\documentclass[a4paper,12pt]{article}
	\usepackage[ansinew]{inputenc}
	\usepackage{ae}	
	\usepackage[brazil]{babel}
	\usepackage{graphicx}
	
	% O pacote psfrag permite a substituição de fontes numa figura EPS por fontes do documento LaTeX (Computer Modern Roman, por exemplo).
	\usepackage{psfrag}
	
	\graphicspath{{Figuras/EPS/}}
	
	\newcommand{\comando}[1]{{\normalfont\textbackslash\color{blue!50!black}#1}}
\newcommand{\pacote}[1]{{\normalfont\sffamily\color{orange!50!black}#1}}
\newcommand{\ambiente}[1]{{\normalfont\sffamily\color{green!50!black}#1}}
\newcommand{\contador}[1]{{\normalfont\sffamily\color{red!50!black}#1}}
\newcommand{\foreign}[1]{\textsl{#1}}
\newcommand\linhabase{\noindent\makebox[0pt][l]{\color{red}\rule{\textwidth}{0.1pt}}}
\begin{document}

	\begin{figure}
		\centering
		\psfrag{Texto na figura}[][]{$\displaystyle\Gamma(z) = \int_0^\infty \lambda^{z-1} e^{-\lambda}\, d\lambda$}
		\includegraphics[width=0.5\textwidth]{Tei-Gi.eps}
		\caption{a substituição de texto só funciona com figuras EPS (\foreign{encapsulated PostScript}), o que significa que, ao optar por este recurso, você deve compilar seu documento para PS (\foreign{PostScript}). Mas é possível produzir um documento PDF através do perfil de compilação \LaTeX\ => PS => PDF.}
	\end{figure}

\end{document}
